\begin{frame}{우리 팀 발표를 리뷰하기: 발표 전 셀프-리뷰}
  \begin{itemize}
    \item 최선의 전략: \textbf{네가 심사할 때 던질 수 있는 가장
      날카로운 질문을, 네 발표 \emph{전에} 스스로 던져보고 답을
      준비해두는 것}.
    \item §5가지 패턴의 5개 질문을 자기 팀 보고서에 적용:
    \begin{itemize}
      \item (1) 시드 3개 greedy 평가 \textbf{표}가 보고서에 있는가?
      \item (2) \textbf{평가 프로토콜 문장}($\varepsilon=0$, 200
        에피소드, 모든 알고리즘·시드에 동일)이 있는가?
      \item (3) 무작위 정책 리턴이 \textbf{실측 숫자}로 적혀 있는가?
      \item (4) $\alpha$·$\varepsilon$ 민감성을 \textbf{한 줄이라도}
        검토했는가?
      \item (5) 불안정한 시드·구간이 ``잘 안 된 부분'' 절에
        \textbf{사라지지 않고} 있는가?
    \end{itemize}
  \end{itemize}
  \vskip0.5em
  \begin{block}{왜 발표 \emph{전}인가}
    이 5가지에 답이 없다면, 발표 전에 해결해두는 것이 발표 중에
    3점 리뷰를 받는 것보다 \textbf{항상 낫다}. Block B Ch\,16.1도
    동일한 구조(발표 $\to$ peer-review) --- 이 셀프-리뷰 습관이
    학기 말까지 이어진다.
  \end{block}
\end{frame}
